Fix an integer $n\ge 2$. Define a decreasing sequence $(a_k)$ of integers in $[1,n)$ by
\[
 a_1:=n-1,
\]
and for $k\ge 2$ define $a_k$ (if it exists) to be the greatest integer $m$ with $1\le m<a_{k-1}$ such that every prime factor $p$ of $m$ satisfies
\[
 p>n-m.
\]
(If no such $m$ exists, the sequence terminates.)

Question: Is it true that for all sufficiently large $n$, the sequence $(a_k)$ cannot consist solely of primes (aside from the terminal $1$)?
Equivalently: for all sufficiently large $n$, does there exist a composite $m$ with $1<m<n$ whose prime factors all exceed $n-m$?

Conventions: $1$ has no prime factors, so it vacuously satisfies the “all prime factors” condition.
